% !TeX program = xelatex
% 每日阅读包 —— A 部分（由 Homework/_review/tools/build_daily_pack.py 生成）
\documentclass[10pt,a4paper]{ctexart}
\usepackage[a4paper,top=1.9cm,bottom=1.9cm,left=1.9cm,right=1.9cm]{geometry}
\usepackage{amsmath,amssymb,amsthm,mathtools,bm}
\usepackage{booktabs,array,longtable,tabularx,multicol}
\usepackage{enumitem}
\usepackage[most]{tcolorbox}
\usepackage{fancyhdr}
\usepackage{titlesec}
\usepackage{hyperref}
\hypersetup{hidelinks,colorlinks=true,linkcolor=blue!50!black}
\linespread{1.20}
\setlength{\parskip}{0.18em}
\renewcommand{\arraystretch}{1.30}
\setlist[enumerate]{itemsep=1pt,topsep=3pt,parsep=0pt,leftmargin=1.7em}
\setlist[itemize]{itemsep=1pt,topsep=3pt,parsep=0pt,leftmargin=1.5em}
\titleformat{\section}{\large\bfseries\color{blue!35!black}}{\thesection}{0.5em}{}
\titleformat{\subsection}{\normalsize\bfseries\color{blue!25!black}}{\thesubsection}{0.45em}{}
\titlespacing*{\section}{0pt}{1.1ex plus .2ex}{0.6ex}
\pagestyle{fancy}
\fancyhf{}
\fancyhead[L]{\small 每日阅读包 · 2026-11-07}
\fancyhead[R]{\small 陶理}
\fancyfoot[C]{\small \thepage}
\renewcommand{\headrulewidth}{0.4pt}
\newcommand{\R}{\mathbb{R}}
\newcommand{\eps}{\varepsilon}
\newcommand{\abs}[1]{\left\lvert #1\right\rvert}
\newcommand{\ee}{\mathrm{e}}
\newcommand{\D}{\mathrm{d}}
\newcommand{\blank}{\underline{\hspace{1.5cm}}}
\newtcolorbox{keypoint}[1][]{enhanced,breakable,colback=blue!3,colframe=blue!40!black,
  fonttitle=\bfseries,title={#1},boxrule=0.5pt,arc=1.2mm,
  left=2.5mm,right=2.5mm,top=1.5mm,bottom=1.5mm,before skip=3pt,after skip=4pt}
\newtcolorbox{warnbox}[1][]{enhanced,breakable,colback=red!2,colframe=red!45!black,
  fonttitle=\bfseries,title={#1},boxrule=0.5pt,arc=1.2mm,
  left=2.5mm,right=2.5mm,top=1.5mm,bottom=1.5mm,before skip=3pt,after skip=4pt}
\begin{document}

\begin{center}
{\Huge\bfseries 每日阅读包}\\[0.25em]
{\LARGE 2026-11-07（周六）}\\[0.45em]
{\Large\bfseries 补齐 11/6（D-4）＋ 总体复习}\\[0.3em]
{\normalsize 中值定理保底 · 加练 50 题 · 高难 22 题冲顶 · 九大专题快扫}
\end{center}
\vspace{0.4em}
\begin{keypoint}[今天的一句话]
\textbf{不做整卷}（2024-25 真题卷之前已做过）。今天只有一组「新内容」——补齐 11/6 的
中值定理；其余全部是\textbf{闭卷默写 + 动笔做题 + 订正}。
\textbf{15:23 改订}：数列极限 / 新压轴已按你的要求压缩为最小保温（B6 只 2 页、S 组只 3 题），
省下的时间给 \textbf{B7--B8《今日讲义配套加练》}（50 题题面 + 逐题详解）。
\textbf{16:30 追加}：新编 \textbf{B14--B15《高难冲顶 22 题》}——对标 2008--2025 真题的中高难度与压轴，
题面与详解分开；\textbf{这是冲高块，不替代 P0 的 11/6 补齐}。
只选一样时永远选「动手做题」，\textbf{不通读任何一本笔记}。
\end{keypoint}
\section*{五级优先级（时间不够从下往上砍）}
\begin{tabularx}{\linewidth}{c|c|p{8.1cm}|c}
\toprule
级 & 时长 & 内容 & 必做 \\
\midrule
P0 & 150 min & 补齐 11/6：中值定理保底 + 加练讲义 M 组 8 题 + 1 页证明模板卡 & \textbf{必做} \\
P1 & 115 min & 总体复习：九大专题快扫，每格「闭卷默写 + 1 道代表题」直接取 B7 对应组 & \textbf{必做} \\
P2 & 90 min & 散项收口：Q 卷 30 min 掐表 + 订正 / §7 得分补记 / *加练 / L 组 / 两道 Q17 & 按余力 \\
\textbf{P2+} & 100 min & \textbf{冲高块：高难 22 题（B14--B15）}。整册 110 min；余力不足只做六道优先题（约 50 min） & 按余力 \\
P3 & 15 min & 睡前只回忆 5 张手抄卡 & 固定 \\
\bottomrule
\end{tabularx}
\section*{这份 PDF 怎么用}
\begin{enumerate}
\item \textbf{第 2 页}是阅读路线图：B1--B15 对应后面合订的材料原件，按段号顺序翻，不要跳。
\item \textbf{第 3--6 页}是\textbf{可书写页}：闭卷默写区、记录表、验收清单。建议打印，或平板手写。
\item \textbf{核对答案}放在默写页之后：\textbf{做完再看}，否则闭卷无效。
\item \textbf{B7--B8 是新编配套加练讲义}：B7 是 50 道题面（M/L/J/T/G/E/S 七组 + Q 卷），
      B8 是答案速查、逐题详解、错因回炉索引与验收清单。\textbf{按题组与 B2--B5、B9--B13 交叉使用}，
      不必一次做完。
\item \textbf{B14--B15 是 16:30 追加的高难册}：B14 是 22 道题面 + 难度对标表，B15 是答案与逐题详解。\textbf{每题先掐 8--12 min；卡 5 min 只看「卡点」一行，合上重写}。时间不足按 B14 里列的六道优先题做（约 50 min），\textbf{不要为了做完而跳过 P0}。
\item \textbf{第 11 页起}是今天要看的材料原件，每页左下角标了来源与原页码。
      PDF 书签（大纲）可按 B1--B15 直接跳转。
\item 今天\textbf{不动 11/8 起的排期}；Level 4--6（积分 / 多元 / 级数与 ODE）不在本次期中范围。
\end{enumerate}
\begin{warnbox}[产出（明天 11/8 的输入）]
(1) 2024-25 旧卷得分与错题号补记；(2) 1 页证明模板卡；(3) 九宫格快扫记录；(4) 3 条最大风险；
(5) 加练讲义 Q 卷得分（$\geq 14/20$ 达标）；(6) \textbf{高难 22 题完成题数}（能独立写出第一行的题数，目标 $\geq 16$）。
\textbf{旧卷得分全库无记录，记不清就写「未记录」，不要估算分数。}
\end{warnbox}

\clearpage

\section*{阅读路线图（按段号顺序翻）}
{\small 表中「页」是\textbf{本 PDF 的绝对页码}；原件页码写在材料列的 p 后面。}
\small
\setlength{\tabcolsep}{3pt}
\begin{longtable}{c|c|p{4.5cm}|p{6.6cm}|c}
\toprule
段 & 页 & 材料（原件页码） & 翻到这里做什么 & min \\
\midrule
\endhead
B1 & 11 & 讲义\_该翻哪一本.pdf p2,3 & 开局读总规则；P1 九宫格按此表逐格走 & 5 \\
B2 & 13 & 微分中值定理专题讲义.pdf p9,10,11,12 & P0(3) 遮住解答，重推 6 个辅助函数 + 三道倒推示范 & 35 \\
B3 & 17 & 微分中值定理专题讲义.pdf p13,14,15 & P0(4) 挑 3 类：$f'+\lambda f=0$ / 含 $f''$ 两次 Rolle / 结论含 $\ln$ 的 Cauchy & 25 \\
B4 & 20 & 微分中值定理专题讲义.pdf p16,17,18 & P0(5) 做 2010 级或 2017 级 七(18)，按答题模板五步写全 & 20 \\
B5 & 23 & 高数A类\_第1学期\_期中\_2009-2010学年.pdf p1,2 & P0(5) 做 Q18（$f(c)f'(1-c)$ 型构造），条件与端点验证写全 & 15 \\
B6 & 25 & 讲义\_数列极限与压轴新题型.pdf p1,4 & 只闭卷复述 P1--P3 三问思路；动笔只做加练讲义 S 组 3 题（B7）。原 §3 十例 / V1--V5 / 模拟卷 C Q18 今天不做 & 8 \\
B7 & 27 & 讲义\_今日配套加练\_11-07.pdf p1,2,3,4,5 & 今天的动笔主战场：P0 后插 M 组 40 min；P1 各格「1 道代表题」直接改用本册对应组；Q 卷 30 min 掐表 & 70 \\
B8 & 32 & 讲义\_今日配套加练\_11-07.pdf p6,7,8,9 & 做完再看；错题按第 12 节回炉索引定位，错的「第一步」写进记录表 C & 30 \\
B9 & 36 & 讲义\_极限定义证明.pdf p4,5,6,14 & P1(1) 三形式框架各写一遍，重点 M 型 $\delta=\min\{1/2,1/(2G)\}$ & 15 \\
B10 & 40 & 讲义\_函数极限与间断点.pdf p3,4,7,8,11,12,13,14 & P1(2)(3) 挑 1 道有理化 + 1 道分界点；§7 得分补记；§9 全表打勾 & 27 \\
B11 & 48 & 讲义\_练习\_函数极限计算与泰勒.pdf p3,5,6,7,8 & P1(4) B 组挑 1 题；P2(3) L 组 30 min 掐表，目标 $\geq 14/20$ & 42 \\
B12 & 53 & 讲义\_导数计算专题.pdf p2,3,4,5 & P1(5) 挑 1 道分段可导或参数二阶导；§7 速查卡 & 15 \\
B13 & 57 & 讲义\_函数性态与作图.pdf p2,3,4 & P1(6) 挑 1 道例题只列七步表；P2(4) 两道 Q17 闭卷七步 & 35 \\
B14 & 60 & 讲义\_高难冲顶\_11-07.pdf p1,2,3,4 & P2+ 冲高块：整册 110 min；余力不足只做 H2.9 / H2.10 / H3.11 / H3.13 / H4.16 / H5.20 六道 & 70 \\
B15 & 64 & 讲义\_高难冲顶\_11-07.pdf p5,6,7,8,9,10,11,12,13,14,15 & 做完一组才翻；卡 5 min 只看「卡点」一行，合上重写 & 40 \\
\bottomrule
\end{longtable}
\normalsize
\begin{keypoint}[路线图的两条纪律]
(1) 每一段都是「\textbf{先遮住解答自己做}，再对答案」；看不懂才回笔记查一步，\textbf{不从头读笔记}。
(2) B7 的题面与 B8 的详解\textbf{必须分开用}：做完一组才翻 B8 对答案。
(3) 数列极限（B6、九宫格第 7 格、S 组）已压缩，\textbf{只保温不展开}；中值定理两格在 P1 中与 P0 合并计时，\textbf{不重复做}；整卷一格今天只口述 120 分钟分配表。
(4) \textbf{B14 与 B15 必须分开用}：B14 是 22 道题面，B15 是答案与详解。做完一组才翻 B15；卡住只看详解里的「卡点」一行，\textbf{不要顺着读下去}。
\end{keypoint}

\clearpage

\section*{P0 闭卷默写（15 min，写不出的标「?」，不翻书）}
\textbf{(1) 三定理的条件与结论：} Rolle / Lagrange / Cauchy，各写「闭区间 / 开区间 / 端点」三行。
\par\smallskip\noindent\rule{\linewidth}{0.25pt}\par\vspace{0.52cm}
\noindent\rule{\linewidth}{0.25pt}\par\vspace{0.52cm}
\noindent\rule{\linewidth}{0.25pt}\par\vspace{0.52cm}
\noindent\rule{\linewidth}{0.25pt}\par\vspace{0.52cm}
\textbf{(2) 六个辅助函数（写目标结论 $\to$ 构造 $F(x)$）：}
\par\smallskip\noindent\rule{\linewidth}{0.25pt}\par\vspace{0.52cm}
\noindent\rule{\linewidth}{0.25pt}\par\vspace{0.52cm}
\noindent\rule{\linewidth}{0.25pt}\par\vspace{0.52cm}
\noindent\rule{\linewidth}{0.25pt}\par\vspace{0.52cm}
\noindent\rule{\linewidth}{0.25pt}\par\vspace{0.52cm}
\noindent\rule{\linewidth}{0.25pt}\par\vspace{0.52cm}
\textbf{(3) 两个中值点模型、Rolle 答题五步、条件三查、证不等式三条路：}
\par\smallskip\noindent\rule{\linewidth}{0.25pt}\par\vspace{0.52cm}
\noindent\rule{\linewidth}{0.25pt}\par\vspace{0.52cm}
\noindent\rule{\linewidth}{0.25pt}\par\vspace{0.52cm}
\noindent\rule{\linewidth}{0.25pt}\par\vspace{0.52cm}
\noindent\rule{\linewidth}{0.25pt}\par\vspace{0.52cm}
\noindent\rule{\linewidth}{0.25pt}\par\vspace{0.52cm}
\noindent\rule{\linewidth}{0.25pt}\par\vspace{0.52cm}
\begin{keypoint}[验收]
6 个辅助函数闭卷全对；答题模板五步不看资料写出；能说清「Peano 余项不能证严格不等式」。
\end{keypoint}

\clearpage

\section*{P1 九宫格快扫：闭卷默写（每格 2 min）}
{\small 统一动作：2 min 闭卷默写决策卡 $\to$ 8--12 min 只做 1 道代表题 $\to$ 2 min 对答案写一句错因。}\par
\small
\setlength{\tabcolsep}{4pt}
\renewcommand{\arraystretch}{2.45}
\par
\begin{center}
\begin{tabular}{p{3.0cm}|p{13.4cm}}
\toprule
专题 & 闭卷默写（写关键词，不抄书） \\
\midrule
(1) 第 11 题定义证明 & \hrulefill \\
(2) 函数极限计算 / 三型 & \hrulefill \\
(3) 间断点 & \hrulefill \\
(4) 泰勒 & \hrulefill \\
(5) 导数计算 & \hrulefill \\
(6) 函数性态作图 & \hrulefill \\
(7) 数列极限（\textbf{已压缩}，只默写四法名） & \hrulefill \\
(8) 中值定理（与 P0 合并） & \hrulefill \\
(9) 高频四件 + 时间分配 & \hrulefill \\
\bottomrule
\end{tabular}
\end{center}
\renewcommand{\arraystretch}{1.30}
\normalsize
\vspace{0.4em}
\textbf{动笔题直接取 B7 的 M / L / J / T / G / E 组}：每格做完「闭卷默写」就在对应组挑 1 题，
不再另外找题；做完才翻 B8 对答案。S 组 3 题做完即止，\textbf{不追加}。
\vspace{0.4em}

\textbf{高频必考四件（各 2 min，穿插在快扫里）：}
\par\smallskip\noindent\rule{\linewidth}{0.25pt}\par\vspace{0.52cm}
\noindent\rule{\linewidth}{0.25pt}\par\vspace{0.52cm}
\noindent\rule{\linewidth}{0.25pt}\par\vspace{0.52cm}
\textbf{每格错的「第一步」汇总成 3 条最大风险：}
\par\smallskip\noindent\rule{\linewidth}{0.25pt}\par\vspace{0.52cm}
\noindent\rule{\linewidth}{0.25pt}\par\vspace{0.52cm}
\noindent\rule{\linewidth}{0.25pt}\par\vspace{0.52cm}
\clearpage

\section*{记录表 A：2024-25 旧卷补记（11/8 对比表唯一输入）}
\begin{warnbox}[先补记，再开始 P0]
全库查不到那次整卷的得分与错题号。\textbf{记不清写「未记录」，不要估算分数}。
\end{warnbox}
\setlength{\tabcolsep}{5pt}
\renewcommand{\arraystretch}{2.35}
\begin{tabular}{p{5.2cm}|p{11.1cm}}
\toprule
项 & 填写 \\
\midrule
那次得分 & \hrulefill \\
失分题号 & \hrulefill \\
三个最大失分点 & \hrulefill \\
最耗时三题 & \hrulefill \\
\bottomrule
\end{tabular}
\renewcommand{\arraystretch}{1.30}
\vspace{0.8em}
\section*{记录表 B：九宫格快扫（每格一行）}
\small\setlength{\tabcolsep}{3pt}
\renewcommand{\arraystretch}{2.10}
\begin{tabularx}{\linewidth}{c|p{3.0cm}|p{3.5cm}|X}
\toprule
格 & 默写了什么 & 做的哪道题 & 错的第一步 / 触发条件 \\
\midrule
1 & & & \\
2 & & & \\
3 & & & \\
4 & & & \\
5 & & & \\
6 & & & \\
7 & & & \\
8 & & & \\
9 & & & \\
\bottomrule
\end{tabularx}
\renewcommand{\arraystretch}{1.30}
\normalsize

\clearpage

\section*{记录表 C：错因 $\to$ 触发条件（每题一句）}
{\small 格式：「下次看到 \blank 先做 \blank」。只写第一步，不抄整题。}
\par\smallskip\noindent\rule{\linewidth}{0.25pt}\par\vspace{0.58cm}
\noindent\rule{\linewidth}{0.25pt}\par\vspace{0.58cm}
\noindent\rule{\linewidth}{0.25pt}\par\vspace{0.58cm}
\noindent\rule{\linewidth}{0.25pt}\par\vspace{0.58cm}
\noindent\rule{\linewidth}{0.25pt}\par\vspace{0.58cm}
\noindent\rule{\linewidth}{0.25pt}\par\vspace{0.58cm}
\noindent\rule{\linewidth}{0.25pt}\par\vspace{0.58cm}
\noindent\rule{\linewidth}{0.25pt}\par\vspace{0.58cm}
\section*{记录表 D：3 条最大风险 + 明天要防的复发点}
\par\smallskip\noindent\rule{\linewidth}{0.25pt}\par\vspace{0.58cm}
\noindent\rule{\linewidth}{0.25pt}\par\vspace{0.58cm}
\noindent\rule{\linewidth}{0.25pt}\par\vspace{0.58cm}
\noindent\rule{\linewidth}{0.25pt}\par\vspace{0.58cm}
\noindent\rule{\linewidth}{0.25pt}\par\vspace{0.58cm}
\section*{今日验收清单}
\begin{multicols}{2}
\begin{itemize}
\item[$\square$] 旧卷得分 / 错题已补记（或明确写「未记录」）
\item[$\square$] 6 个辅助函数闭卷全对；Rolle 五步不看资料写出
\item[$\square$] 2009-10 Q18 + 一道 2010/2017 级七(18) 按模板写全
\item[$\square$] 加练讲义 M / L / J / T 四组各至少完成一半，S 组 3 题做完即止
\item[$\square$] Q 卷 30 分钟掐表，$\geq 14/20$ 达标
\item[$\square$] 错的「第一步」写进记录表 C
\item[$\square$] 九大专题每格：闭卷默写 + 1 道代表题（B7 组）
\item[$\square$] \textbf{高难 22 题}：六道优先题（H2.9 / H2.10 / H3.11 / H3.13 / H4.16 / H5.20）至少做完 4 道
\item[$\square$] 压轴三道（H3.11 / H3.13 / H5.22）能不看提示写出第一行
\item[$\square$] §7 得分补记、*加练、L 组复测至少各完成一半
\item[$\square$] 手抄 5 张卡齐：定义三模板 / 泰勒决策 / 求导检查 / 作图七步 / 证明模板
\item[$\square$] 23:30 前睡
\end{itemize}
\end{multicols}

\clearpage

\section*{核对答案 (1)：中值定理（做完 P0 再看）}
\textbf{三定理}：Fermat（内点取极值 + 可导 $\Rightarrow f'(x_0)=0$）；
Rolle（$[a,b]$ 连续、$(a,b)$ 可导、$f(a)=f(b)\Rightarrow\exists\xi:f'(\xi)=0$）；
Lagrange（$f(b)-f(a)=f'(\xi)(b-a)$，用于证不等式与估计函数差）；
Cauchy（$\frac{f(b)-f(a)}{g(b)-g(a)}=\frac{f'(\xi)}{g'(\xi)}$，L'Hôpital 由它推出）。
\small\setlength{\tabcolsep}{4pt}
\begin{tabularx}{\linewidth}{p{4.6cm}|p{4.4cm}|X}
\toprule
目标结论 & 构造 $F(x)$ & 对谁用 Rolle \\
\midrule
$f'(\xi)+\lambda f(\xi)=0$ & $\ee^{\lambda x}f(x)$ & 两个零点 /
$F(a)=F(b)$ \\
$\xi f'(\xi)+n f(\xi)=0$ & $x^n f(x)$（$0\notin[a,b]$） & 同上 \\
$\xi f'(\xi)-n f(\xi)=0$ & $f(x)/x^n$ & 同上 \\
$f'(\xi)+g'(\xi)f(\xi)=0$ & $\ee^{g(x)}f(x)$ & 指数因子一般化 \\
$f'(\xi)=k$ & $f(x)-kx$ & 退化为 Rolle \\
$f'(\xi)=2\xi$ 一类 & $f(x)-x^2$ & 右端积分回去 \\
\bottomrule
\end{tabularx}
\normalsize\smallskip
\textbf{倒推口令}：把要证的式子看成「某个乘积的导数为零」；$u'+\frac{u'}u u$ 型想 $\ln\abs u$
或 $\ee^{\lambda x}$；$\xi f'+nf$ 就是 $(x^n f)'$ 除以 $x^{n-1}$。
\textbf{Rolle 五步}：造 $F$ $\to$ 验连续 / 可导 $\to$ 验 $F(a)=F(b)$ $\to$ 写定理得 $\xi$ $\to$
回代化简。
\textbf{条件三查}：闭区间连续？开区间可导？端点值相等 / $g'\ne0$？缺一条结论就可能假。
\textbf{两个中值点}：分区间两段各用一次，或对不同函数各用一次。
\textbf{证不等式三条路}：单调性（移项造 $g$，$g(a)=0$ 且 $g'>0$）、Lagrange 放缩、
Taylor 带 \textbf{Lagrange} 余项（Peano 余项不能给严格不等号）。

\section*{核对答案 (1)（压缩版）：数列极限只保温到这一行}
四法只记名字与触发条件：\textbf{夹逼}（两端同极限、放缩分母取最值）、
\textbf{单调有界}（先证单调 + 有界再对递推取极限）、
\textbf{定积分定义}（认准 $\frac1n\sum f(\frac kn)$）、\textbf{Stolz}（$b_n\nearrow+\infty$）。
压轴三型只记构造：不动点看 $g=f(x)-x$ 两端符号；等距平移把 $n$ 段函数值相加反证；
迭代收敛用 $\abs{a_{n+1}-L}\le q\abs{a_n-L},\ 0<q<1$。

\begin{keypoint}[15:23 改订：这一段今天不展开]
用户明确「最后数列极限那部分可以减少」：B6 由 7 页压到 2 页（只留讲义 p1 母题与 p4 速查卡），
\textbf{删去} §3 十道例题、V1--V5 限时卷、模拟卷 C 第 18 题换题。
动笔只做 B7 加练讲义 \textbf{S 组 3 题}（定积分定义 / 递推不动点 / 导数定义型），
2024-25 Q18 已随旧卷做过，今天\textbf{连口述都省了}。省下的时间全部给 M / L / J / T / G / E 六组与 Q 卷。
完整四法与 P1--P3 母题仍在 \texttt{讲义\_数列极限与压轴新题型.pdf}，考后复盘再翻。
\end{keypoint}

\clearpage

\section*{核对答案 (2)：导数计算六型}
\textbf{动笔前三问}：定义域？前提（$x'_t\ne0$ / $F_y\ne0$ / 分段点先验连续）？用哪条公式？
\small\setlength{\tabcolsep}{4pt}
\begin{tabularx}{\linewidth}{c|p{3.4cm}|p{5.2cm}|X}
\toprule
型 & 触发 & 第一动作 & 必写前提 \\
\midrule
D1 & 分段可导求参数 & 先列连续方程 $f(a^-)=f(a^+)$，再列导数方程 $f'_-(a)=f'_+(a)$ & 分界点两侧用定义 \\
D2 & 参数方程二阶导 & 先算 $y'/x'$，再对 $t$ 求导、除以 $x'$（\textbf{不是} $y''/x''$） & $x'(t)\ne0$ \\
D3 & 隐函数二阶导 & $y'=-F_x/F_y$，再对 $x$ 求导并把 $y'$ 代回 & $F_y\ne0$ \\
D4 & Leibniz / 高阶导 & 找「低阶导归零」的因子，只留 $\binom nk$ 中那几项 & 展开式写到 $n$ \\
D5 & 极坐标切线 & 先化 $(x,y)=(r\cos\theta,r\sin\theta)$，斜率 $\frac{r'\sin\theta+r\cos\theta}{r'\cos\theta-r\sin\theta}$ & $\theta$ 范围、$r<0$ 反向 \\
D6 & 微分 $\D y$ & 求导后乘 $\D x$ & 一阶微分形式不变性 \\
\bottomrule
\end{tabularx}
\normalsize\smallskip
\textbf{达布定理}：区间上处处可导 $\Rightarrow f'$ 无第一类间断点（\textbf{只在处处可导时}；「左 $\le$ 右」不是定理）。

\section*{核对答案 (2)：函数性态作图七步}
\begin{enumerate}
\item 定义域（\textbf{含不可导点}）$\to$ (2) 奇偶 / 周期 $\to$ (3) 渐近线（\textbf{两端分别求}）$\to$
      (4) $f'$ 判单调与极值 $\to$ (5) $f''$ 判凹凸与拐点 $\to$ (6) \textbf{列表} $\to$
      (7) 作图（标极值点、拐点、交点、渐近线）。
\end{enumerate}
\textbf{渐近线三查}：垂直＝无定义点看单侧 $\lim f=\infty$；水平＝$\lim_{x\to\pm\infty}f=b$；
斜＝$k=\lim f/x,\ b=\lim(f-kx)$。两端各查一次；二次增长（$f/x\to\infty$）判「无斜」。
\textbf{列表是给分点}：区间 / $f'$ 符号 / 升降 / $f''$ 符号 / 凹凸，五行齐全；
卷面附了 $y',y''$ 就从第 (4) 步开始，别回头重算。不可导点在第 (1) 步单列（$f'=\infty$ 是竖直切线，不是驻点）。

\clearpage

\section*{核对答案 (3)：计算题三型（限时卷 9 / 10 / 11）}
\small\setlength{\tabcolsep}{4pt}
\begin{tabularx}{\linewidth}{c|c|X}
\toprule
型 & 答案 & 第一动作（背这句） \\
\midrule
A 根式差 / 加减相消 & $\frac14$ & 首项会相消的差式\textbf{绝不逐项等价} $\to$ 先有理化，把差变成积 \\
B $1^\infty$ 幂指 & $\ee^{1/3}$ & 写 $u^v=\ee^{v\ln u}$，只需把 $u-1$ 展到与 $1/v$ 同阶 \\
C $\lim_n$ 定义的分段函数 & $x=\pm1$ 均跳跃 & 先写三段 $f(x)$，再列可疑点算左右极限 \\
\bottomrule
\end{tabularx}
\normalsize\smallskip
\textbf{$1^\infty$ 捷径} $\lim u^v=\ee^{\lim v(u-1)}$ 成立三条件：$u\to1$；$\lim v(u-1)$ 存在；
$v(u-1)^2\to0$。\textbf{必背零件}：$\sqrt{1+u}-1\sim\frac u2$，$\tan x-\sin x\sim\frac{x^3}{2}$，
$\tan x-x\sim\frac{x^3}{3}$，$1-\cos x\sim\frac{x^2}{2}$。收尾代 $x=10^{-3}$ 扫符号与量级。

\section*{核对答案 (3)：泰勒决策卡 / 间断点 / 定义三模板}
\textbf{泰勒总纲}：先问「是不是 $0/0$ 或 $\infty/\infty$」「分母最高几阶」，$k$ 阶就展到 $k$ 阶；
\textbf{乘积 / 商里的因子}用等价无穷小，\textbf{和差里要消低阶项}必须泰勒（$\sin x-x$ 不能替换 $\sin x\sim x$）；
条件给极限求某阶导数就写 $f(x)=f(0)+f'(0)x+\frac{f''(0)}2x^2+o(x^2)$ 比系数。
注意 $\sqrt{1-x}=1-\frac x2-\frac{x^2}{8}+\cdots$（$x^2$ 项是负号）。
\vspace{0.2em}

\textbf{间断点判型}：先算可疑点左右极限。左右都存在且相等（但无定义或取值不等）$\to$ 可去；
都存在但不等 $\to$ 跳跃（第一类）；一侧 $\pm\infty$ $\to$ 无穷（第二类）；都不存在 $\to$ 振荡（第二类）。
\textbf{八类必分左右}：$\abs x$、$\ee^{1/x}$、$\arctan(1/x)$、$[x]$、$\lim_n$ 定义的分段函数、分母为零点、
分段点、$\sin(1/x)$。2018-19 官方给 1 个第一类间断点是错的，正确 \textbf{2} 个（$x=0$ 处 $\arctan(x+\frac1x)$ 跳跃）。
\vspace{0.2em}

\textbf{定义三模板}：
A（$\eps$-$\delta$）$\forall\eps>0\,\exists\delta>0:0<\abs{x-x_0}<\delta\Rightarrow\abs{f(x)-A}<\eps$；
B（$\eps$-$X$）$\forall\eps>0\,\exists X>0:x>X\Rightarrow\abs{f(x)-A}<\eps$；
C（无穷大 $M$ 型）$\forall M>0\,\exists\delta>0:0<\abs{x-x_0}<\delta\Rightarrow\abs{f(x)}>M$（\textbf{方向相反}）。
收尾：A 取 $\delta=\min\{\delta_1,\eps/C\}$，B 取 $X=\max\{X_1,C/\eps\}$，
C 取 $\delta=\min\{\delta_1,k/M\}$。
\begin{warnbox}[2024-25 第 11 题答案册放缩方向错]
$\frac1{|x|(1+|x|)}\ge\frac1{|x|}$ 不成立；官方 $\delta=\min\{\frac12,\frac1G\}$ 有反例
（$G=10,\ x=0.09631$ 时 $\abs f\approx9.497<10$）。
订正 \textbf{$\delta=\min\{\frac12,\frac1{2G}\}$}，见 \texttt{极限定义证明\_题型专题.md} §3.7。
\end{warnbox}

\clearpage

\section*{高频四件与考场 120 分钟分配}
\textbf{高频四件}：(1) 微分 $\D y=f'(x)\D x$（别漏 $\D x$）；(2) Leibniz 高阶导「代点消项」；
(3) 极坐标 $\theta$ 范围（$r<0$ 时点在反向射线）；(4) 间断点\textbf{分左右}。

\textbf{考场时间分配}（只口述，不掐表）：选填 30 分 / 25 min；第 11 题 8 分 / 10 min；
计算 12--14 / 25 min；证明 15--16 / 20 min；作图 17 / 20 min；压轴 18 / 15 min；检查 5 min。

\section*{四处官方答案错误（对答案时以本库复核结论为准）}
\begin{enumerate}
\item 2015-16 上界；\quad (2) 2024-25 第 11 题不等式 / 放缩方向；
\item 2018-19 间断点个数：官方 1，正确 \textbf{2}；
\item 2023-24 Q15(2)：官方 $+\frac{13!}{3240}$，正确 \textbf{$-\frac{13!}{3240}$}（sympy + mpmath 140 位复核）。
\end{enumerate}
2023-24 选 3 = A、选 4 = D 与复核一致。根目录 \texttt{HANDOFF\_PROMPT.md} 是上午快照，含两处已证伪记录，勿当权威。

\section*{边界与明天}
\begin{itemize}
\item 今天\textbf{不做任何整卷}（2024-25 / 2023-24 / 模拟卷 A/B/C 都不掐表）；模拟卷 C 第 18 题今天也不做，留作 11/9 缓冲日备用。
\item \textbf{数列极限已按 15:23 指示压缩}：B6 只 2 页、S 组只 3 题、九宫格第 7 格只默写四法名；\textbf{不追加}夹逼/Stolz 展开与 P1--P3 母题。
\item \textbf{高难 22 题（B14--B15）是冲高块，不是保底块}：\textbf{不替代} P0 的 11/6 补齐，也不替代 B7--B8 的保底 50 题。P0 没做完就先回 P0；做完 P0 且有余力才开 B14。单题 15 min 写不出第一行就标记跳过。
\item \textbf{不动 11/8 起日期}。若 2023-24 卷也已做过，明天开新对话时说明，再改 11/8 为专项回炉日。
\item Level 4--6（积分 / 多元 / 级数与 ODE）是长期概念底座，\textbf{不在本次期中范围}。
\item 2024-25 原卷是扫描件、文本层全空；若需重看某题请用 PDF 阅读器看图。
\item 数学公式只认 \texttt{.tex} 源或渲染 PNG，\textbf{不要}用 \texttt{pdftotext} / \texttt{get\_text()} 核对公式。
\end{itemize}

\vfill
\begin{center}
{\small\color{gray} 生成：\texttt{Homework/\_review/tools/build\_daily\_pack.py} ·
依据 \texttt{11-7\_补齐D-4与总体复习.md} ·
原件合订见 B1--B15（第 11 页起）}
\end{center}

\end{document}
